\documentclass[11pt,a4paper]{article}
\usepackage{lmodern}
\usepackage[T1]{fontenc}
\usepackage{amsmath,amssymb}
\usepackage{graphicx}
\usepackage{longtable,booktabs,array}
\usepackage{float}
\usepackage[margin=25mm]{geometry}
\usepackage{caption}
\usepackage{hyperref}
\hypersetup{colorlinks=true,linkcolor=blue,urlcolor=blue,citecolor=blue}
\setkeys{Gin}{width=0.82\linewidth,keepaspectratio}
\usepackage{amsthm}
\newtheorem{theorem}{Theorem}
\newtheorem{proposition}{Proposition}
\newtheorem{lemma}{Lemma}
\newtheorem{corollary}{Corollary}
\newtheorem{definition}{Definition}
\newtheorem{remark}{Remark}
\newtheorem{observation}{Observation}
\usepackage{calc}
\newcounter{none}
\providecommand{\real}[1]{#1}
\providecommand{\pandocbounded}[1]{#1}
\providecommand{\tightlist}{\setlength{\itemsep}{0pt}\setlength{\parskip}{0pt}}
\title{On the Connection Between the Conjugate Complex Norm and
Square-Quantity Readout}
\author{Noriaki Kihara, WF System Co., Ltd.~--- DOI:
10.5281/zenodo.21127200}
\date{v0.5 --- July 2026}
\begin{document}
\maketitle
\emph{An observational extension from phase representation to a
multi-dimensional sum-of-squares representation including an invisible
axis (pure algebra / non-physical)}

\section{0. Premises and stance (how to read this
note)}\label{premises-and-stance-how-to-read-this-note}

This note starts \textbf{only from conjugate complex numbers and their
geometric interpretation}. It takes the stance of \textbf{not connecting
to physics} (spacetime, special/general relativity, quantum theory,
metric signature). Accordingly, we declare the following explicitly.

\begin{enumerate}
\def\labelenumi{\arabic{enumi}.}
\tightlist
\item
  The coordinate symbols \(x,\ y,\ z_1,\ z_2,\ z_3,\dots\) that appear
  below are all \textbf{geometric coordinate labels}. \textbf{They are
  not identified with physical spatial or temporal dimensions.}
\item
  In particular, this note \textbf{does not use} a symbol \(t\) for
  time. Even when a coordinate taking imaginary values is introduced, it
  is not ``time'' but merely an \textbf{imaginary-valued coordinate
  label} \(z_h\).
\item
  Even if a sign (a minus) of the sum-of-squares-difference type appears
  in a formula, it \textbf{does not assert a metric signature}. The
  minus is nothing more than an \textbf{algebraic consequence} of
  \(i^2=-1\), arising when a coordinate is placed as imaginary.
\item
  All coordinate axes are treated \textbf{on an equal footing}. Here,
  ``equal footing'' means \textbf{not assigning any particular physical
  role or metric signature from the outset}. Splitting, in a later
  section, into real-valued and imaginary-valued components as an
  algebraic representation (§5) does not violate this equal footing.
\item
  This note asserts no new physical law and no new mathematical theorem.
  It is an \textbf{observational note} that re-reads the known conjugate
  complex norm and sum of squares by placing them within the preceding
  square-quantity-readout series.
\end{enumerate}

This declaration is the design foundation of this note, made to remove
the error of v0.1 (assigning a minus sign to a particular axis by hand,
and using the radius symbol in two senses).

\begin{center}\rule{0.5\linewidth}{0.5pt}\end{center}

\section{Abstract}\label{abstract}

This note is the fourth in the series following linear simplexification
by square-quantity readout, the isolated-peak wave by a half-wavelength
odd-harmonic sum, and the contrast law in two-copy common-relative-phase
superposition.

In the preceding square-quantity-readout paper, reading \(X_i=x_i^2\) in
the positive region showed that the sum-of-squares constraint
\(\sum_i x_i^2=E\) can be read as the linear simplex constraint
\(\sum_i X_i=E\) in square-quantity space. Here, we connect to this
square-quantity readout the most basic norm of conjugate complex
numbers,

\[
Z\bar Z=x^2+y^2 .
\]

For the complex polar form
\(Z=\rho e^{i\theta}=\rho(\cos\theta+i\sin\theta)\), putting
\(x=\rho\cos\theta,\ y=\rho\sin\theta\) gives \(Z=x+iy\) and
\(Z\bar Z=x^2+y^2=\rho^2\). This is a standard fact, and there is no
novelty in the formula itself. What this note observes is that this
known norm connects naturally to the square-quantity readout as an
\textbf{operation that bundles the square quantities of two orthogonal
real components into a single real norm}.

Next, we generalize this to an \textbf{equal-footing multi-dimensional
extension}. Bundling real components two at a time into complex numbers
gives

\[
\sum_j Z_j\bar Z_j=\sum_n x_n^2=R^2 ,
\]

and a coordinate axis not displayed in the complex plane (what we call
an invisible axis) simply joins as an \textbf{equal-footing axis with
the same sign \(+\)}, \(z_1,z_2,\dots\). There is no special axis here.

Finally, to the front-half conjugate norm \(\sum_j Z_j\bar Z_j\)
(Hermitian, positive-definite) we add the square \((i z_k)^2\) of an
imaginary-valued coordinate, and observe the \textbf{zero-sum-of-squares
form}

\[
\sum_j Z_j\bar Z_j+\sum_k (i z_k)^2=0 .
\]

Here the conjugate product \(X\bar X\) and the non-conjugate square
\(X^2\) \textbf{agree for real coordinates but diverge only for
imaginary coordinates} (\(X_h\bar X_h=z_h^2\) versus \(X_h^2=-z_h^2\)),
and this divergence is the sole origin of the minus sign. Expanding
gives \(\sum_j Z_j\bar Z_j=\sum_k z_k^2\), i.e.

\[
\sum_n x_n^2=\sum_k z_k^2 .
\]

The minus sign appearing here is an algebraic consequence of \(i^2=-1\),
not a hand-placed metric signature. The conjugate-norm type \(Z\bar Z\)
and the non-conjugate square \(X_h^2\) of an imaginary coordinate are
not the same operation, and this note keeps the distinction. All of this
is a geometric observation on the \textbf{algebraic arrangement} of
conjugate complex numbers and sums of squares; it makes no claim about
physical spacetime or relativity.

\begin{center}\rule{0.5\linewidth}{0.5pt}\end{center}

\section{1. Position within the
series}\label{position-within-the-series}

This note is not a paper on central-projection geometry itself. The
central-projection series takes the curvature radius, geodesic cells,
and angle/area/volume corrections as its subjects, and deliberately kept
its distance from conjugate-complex representations. Inserting the
content of this note directly into the central-projection series would
shift the center of gravity of the argument.

This note belongs, rather, to the following series.

\begin{enumerate}
\def\labelenumi{\arabic{enumi}.}
\tightlist
\item
  Linear simplexification by square-quantity readout, and a
  dimension-by-dimension organization of curvature-correction
  candidates.
\item
  The isolated-peak wave of a constant-amplitude odd-harmonic sum on a
  half-wavelength phase interval.
\item
  Waveform invariance and the contrast law in two-copy
  common-relative-phase superposition.
\item
  This note, i.e.~the connection between the conjugate complex norm and
  square-quantity readout, and its equal-footing multi-dimensional
  extension.
\end{enumerate}

Following these, this note treats \(Z\bar Z=x^2+y^2\) as the minimal
contact point linking phase representation, conjugate complex numbers,
and square-quantity readout, and generalizes it to an equal-footing
multi-dimensional sum of squares.

\begin{center}\rule{0.5\linewidth}{0.5pt}\end{center}

\section{2. Complex polar form and sum of
squares}\label{complex-polar-form-and-sum-of-squares}

\subsection{2.1 Complex polar form}\label{complex-polar-form}

Write a complex number \(Z\) as
\(Z(\rho,\theta)=\rho(\cos\theta+i\sin\theta)\), where \(\rho\ge0\) is
the radius in the complex plane and \(\theta\) the phase angle. Putting
\(x=\rho\cos\theta,\ y=\rho\sin\theta\) gives \(Z=x+iy\). The conjugate
is \(\bar Z=x-iy\), so

\[
Z\bar Z=(x+iy)(x-iy)=x^2+y^2 .
\]

From the polar form as well,
\(Z\bar Z=\rho e^{i\theta}\rho e^{-i\theta}=\rho^2\). Hence

\[
Z\bar Z=x^2+y^2=\rho^2 .
\]

\subsection{2.2 Interpretation as square-quantity
readout}\label{interpretation-as-square-quantity-readout}

In the preceding square-quantity-readout paper, reading \(X_i=x_i^2\) in
the positive region organized the sum-of-squares constraint
\(\sum_i x_i^2=E\) as the linear constraint \(\sum_i X_i=E\). This
note's \(Z\bar Z=x^2+y^2\) is the \textbf{two-component version} of this
square-quantity readout, in which two square quantities \(x^2,\ y^2\)
are bundled into a single real norm as the conjugate product
\(Z\bar Z\).

\subsection{\texorpdfstring{2.3 Two roles of \(i\) (a distinction made
in this
note)}{2.3 Two roles of i (a distinction made in this note)}}\label{two-roles-of-i-a-distinction-made-in-this-note}

This note clearly distinguishes two different roles of the imaginary
unit \(i\). Neglecting this distinction produces errors in the
multi-dimensional extension below.

\begin{itemize}
\tightlist
\item
  \textbf{Role (1): bundling.} When \(Z=x+iy\) with \(x,y\) both real,
  \(Z\bar Z=x^2+y^2>0\). Here \(i\) is merely a notational device for
  gathering two real coordinates into one complex coordinate. The norm
  is always non-negative.
\item
  \textbf{Role (2): imaginary-valued coordinate label.} When the
  coordinate label itself is specified as an imaginary value, we set
  \(X_h=i\,z_h\) (\(z_h\) a real label). Then
  \(X_h^2=(i z_h)^2=-z_h^2\), producing a \textbf{negative squared
  contribution}.
\end{itemize}

Role (1) gives \(x^2+y^2\) (positive-definite); role (2) gives
\(-z_h^2\) (a negative contribution). The two are distinct and must not
be treated as continuous with each other. More precisely, role (1)
appears in the \textbf{conjugate product} \(Z\bar Z\) (Hermitian,
positive-definite), while role (2) appears in the \textbf{non-conjugate
square} \(X_h^2\) (a complex bilinear form). This difference in
\textbf{presence or absence of conjugation} becomes decisive in §5.

\subsection{2.4 The location of novelty}\label{the-location-of-novelty}

\(Z\bar Z=x^2+y^2\) is a standard fact. This note does not present this
formula as a new theorem. Its aim is to place the known complex norm
within the square-quantity-readout series and make explicit the contact
point

\[
\text{square-quantity readout}\quad\leftrightarrow\quad\text{conjugate complex norm},
\]

and to generalize it to an equal-footing multi-dimensional sum of
squares.

\begin{center}\rule{0.5\linewidth}{0.5pt}\end{center}

\section{3. Reading as a phase
representation}\label{reading-as-a-phase-representation}

In the complex polar form, \(\theta\) is an angle in the plane and can
also be read as the phase angle of a wave. That is,

\[
\operatorname{Re}Z=\rho\cos\theta,\qquad \operatorname{Im}Z=\rho\sin\theta,
\]

and \(Z=\rho e^{i\theta}\) is a representation that expands a state with
phase \(\theta\) into two orthogonal components in the complex plane.

What we treat at this stage is the map
\((\rho,\theta)\longmapsto(x,y)\), from a polar form containing a phase
angle to two orthogonal components. No coordinate of time or propagation
direction is introduced here (as declared in §0, the time symbol \(t\)
is not used).

\begin{center}\rule{0.5\linewidth}{0.5pt}\end{center}

\section{4. Extension to an equal-footing multi-dimensional sum of
squares}\label{extension-to-an-equal-footing-multi-dimensional-sum-of-squares}

\subsection{4.1 Bundling by multiple complex
numbers}\label{bundling-by-multiple-complex-numbers}

Let the real coordinates be \(x_1,x_2,\dots,x_{2m}\) and bundle them two
at a time into complex numbers,

\[
Z_j=x_{2j-1}+i\,x_{2j},\qquad Z_j\bar Z_j=x_{2j-1}^2+x_{2j}^2 .
\]

Their sum is

\[
\sum_{j=1}^{m} Z_j\bar Z_j=\sum_{n=1}^{2m} x_n^2 .
\]

For an odd count, add the unbundled real coordinate \(x_{2m+1}\) on its
own. In either case,

\[
\sum_n x_n^2=R^2
\]

is a positive-definite sum of squares containing all coordinates on an
equal footing.

\subsection{4.2 The invisible axis is not a special
axis}\label{the-invisible-axis-is-not-a-special-axis}

When a coordinate does not appear in the complex plane currently
displayed (say the \(Z_1\) plane), we call it an \textbf{invisible
axis}. But this means only that it is not carried on that plane in the
display; its \textbf{sign and role are the same as the other axes}.
Adding invisible axes \(z_1,z_2,\dots\) merely amounts to

\[
\sum_n x_n^2 + z_1^2 + z_2^2 + \cdots = R^2 ,
\]

adding terms of \textbf{the same sign \(+\) to the left-hand side so
that \(R\) grows}. The radius \(R\) remains, throughout, the single
``square root of the sum of squares of all coordinates,'' and is never
defined in two ways.

This is the point that removes the v0.1 error. Taking a particular axis
alone and placing it on the other side, or giving it a minus sign,
breaks \(\sum_n x_n^2=R^2\) and makes \(R\) carry two meanings. As long
as the axes are treated equally, that breakdown does not occur.

Note that ``invisible'' here is limited to the \textbf{display-level
meaning} of not appearing in the chosen complex-plane display. It does
not mean a physically unobservable axis, a hidden physical dimension, or
an unknown degree of freedom.

\subsection{4.3 Representation by a tilt angle (spherical
coordinates)}\label{representation-by-a-tilt-angle-spherical-coordinates}

For three coordinates, using the radius \(R\), a tilt angle \(\phi\),
and the phase angle \(\theta\),

\[
x=R\cos\phi\cos\theta,\qquad y=R\cos\phi\sin\theta,\qquad z_1=R\sin\phi ,
\]

we immediately obtain

\[
x^2+y^2+z_1^2
=R^2\cos^2\phi(\cos^2\theta+\sin^2\theta)+R^2\sin^2\phi
=R^2 .
\]

Here \(\phi\) is a tilt angle leaving the complex plane, not an
additional in-plane rotation phase. The two must be distinguished: an
additional in-plane phase \(\alpha\) merely gives
\(x=\rho\cos(\theta+\alpha),\ y=\rho\sin(\theta+\alpha)\), generates no
\(z_1\) component, and leaves \(x^2+y^2=\rho^2\).

\begin{center}\rule{0.5\linewidth}{0.5pt}\end{center}

\section{5. Conjugate-norm-type square quantity and the
zero-sum-of-squares by an imaginary-valued
coordinate}\label{conjugate-norm-type-square-quantity-and-the-zero-sum-of-squares-by-an-imaginary-valued-coordinate}

\subsection{5.1 The difference between the conjugate product and the
square ------ agreement for real coordinates, divergence for imaginary
ones}\label{the-difference-between-the-conjugate-product-and-the-square-agreement-for-real-coordinates-divergence-for-imaginary-ones}

The previous \(\sum_j Z_j\bar Z_j=\sum_n x_n^2=R^2\) is a
\textbf{positive-definite norm containing conjugation} (the Hermitian
form \(\sum_n X_n\bar X_n=\sum_n|X_n|^2\)), which is zero only when all
components are zero.

Here we clarify the relation between the conjugate product \(X\bar X\)
and the non-conjugate square \(X^2\). \textbf{For real-valued
coordinates the two agree}: if \(x\) is real, \(x\bar x=x^2\).
\textbf{They diverge only for imaginary-valued coordinates}: for
\(X_h=i\,z_h\) (\(z_h\) real),

\[
X_h\bar X_h=(i z_h)(\overline{i z_h})=(i z_h)(-i z_h)=z_h^2\ (>0),
\qquad
X_h^2=(i z_h)^2=-z_h^2\ (<0).
\]

That is, \textbf{a negative squared contribution arises ``only'' when
one takes the non-conjugate square of an imaginary-valued coordinate}.
This is the sole origin of the minus signs appearing hereafter, and is
not a hand-placed metric signature. Accordingly, this note \textbf{does
not confuse} the conjugate-norm type \(Z\bar Z\) with the non-conjugate
square \(X_h^2\) of an imaginary coordinate.

\subsection{5.2 The zero-sum-of-squares form (conjugate-norm backbone +
imaginary-coordinate
square)}\label{the-zero-sum-of-squares-form-conjugate-norm-backbone-imaginary-coordinate-square}

To the positive-definite square quantity \(\sum_j Z_j\bar Z_j\) obtained
from the conjugate norm, we add the square \((i z_k)^2\) of
imaginary-valued coordinates, and observe the
\textbf{zero-sum-of-squares form}

\[
\sum_j Z_j\bar Z_j+\sum_k (i z_k)^2=0,
\qquad Z_j=x_{2j-1}+i x_{2j},\ \ z_k\in\mathbb{R}.
\]

Here the visible part is the conjugate norm (Hermitian,
positive-definite) and the invisible part is the square of
imaginary-valued coordinates (non-conjugate); the two are different
operations. Expanding the left-hand side, \((i z_k)^2=-z_k^2\) gives

\[
\sum_j Z_j\bar Z_j-\sum_k z_k^2=0
\quad\Longrightarrow\quad
\sum_j Z_j\bar Z_j=\sum_k z_k^2 .
\]

Since further \(\sum_j Z_j\bar Z_j=\sum_n x_n^2\),

\[
\sum_n x_n^2=\sum_k z_k^2 .
\]

Note that, since the visible coordinates are real, by §5.1 we have
\(x_n\bar x_n=x_n^2\). Hence this form coincides numerically with the
non-conjugate quadratic form written as \(\sum_n X_n^2=0\) for all
coordinates at once. However, this note takes as primary the above form,
written by \textbf{separating} the visible part as a conjugate norm and
the invisible part as an imaginary-coordinate square, in order not to
confuse the \(Z\bar Z\) type (conjugate) with the \(X^2\) type
(non-conjugate).

\subsection{5.3 A non-trivial solution requires a non-real
component}\label{a-non-trivial-solution-requires-a-non-real-component}

If all coordinates are real, \(\sum_j Z_j\bar Z_j=\sum_n x_n^2=0\) has
only the all-zero solution (a positive-definite norm). Therefore, for
the zero-sum-of-squares to be \textbf{non-trivial, at least one non-real
component is required}. As the minimal organization, this note treats
the non-real component only as a purely imaginary-valued coordinate
\(i z_k\) (in general the form \(a_n+i b_n\) is also possible). The
non-real component is not an additional assumption but appears as a
necessary condition for non-triviality.

\subsection{5.4 The origin of the minus sign, and
non-identification}\label{the-origin-of-the-minus-sign-and-non-identification}

The minus sign appearing here is, as in §5.1, an algebraic consequence
of \(i^2=-1\) from the square of an imaginary-valued coordinate, not a
metric signature assigned to a particular axis by hand. Moreover, this
note does not identify \(z_k\) with a time coordinate, a curvature
radius, or a projection radius (§0, §6).

\subsection{5.5 Contrast with the v0.1
error}\label{contrast-with-the-v0.1-error}

v0.1, to the positive-definite three-dimensional sum
\(x^2+y^2+z^2=r^2\), added by hand a square \(t^2\) of a separate symbol
to the right-hand side, writing \(x^2+y^2+z^2=r^2+t^2\). But this was
not an equal-footing dimensional extension (adding same-sign terms to
the left-hand side); it mixed \textbf{two independent operations}: (i)
introducing a new coordinate, and (ii) placing it with a minus sign. As
a result \(\sum_n x_n^2=R^2\) broke down and \(r\) carried two meanings.

In the present form, the minus sign is \textbf{derived} from the
imaginary-coordinate square (\(i^2\)), the visible-part conjugate norm
\(Z\bar Z\) (positive-definite) and the invisible-part
imaginary-coordinate square \(X_h^2\) are clearly separated, and \(R\)
is not used in two ways, so this breakdown does not occur.

\begin{center}\rule{0.5\linewidth}{0.5pt}\end{center}

\section{6. Readings this note does not bring
about}\label{readings-this-note-does-not-bring-about}

We restate the §0 declaration in terms of concrete formulas.

\begin{itemize}
\tightlist
\item
  This note does not introduce a time coordinate. The \(z_k\) on the
  right-hand side of \(\sum_a x_a^2=\sum_k z_k^2\) is an
  imaginary-valued coordinate label, not time.
\item
  This note does not assert a metric signature \((+,+,\dots,-)\). The
  negative squared contribution is algebra arising from \(i^2=-1\), not
  a spacetime metric.
\item
  This note derives none of the light cone, Minkowski space, special
  relativity, or general relativity. Even if
  \(\sum_a x_a^2=\sum_k z_k^2\) formally coincides with those quadratic
  forms, this note makes no such identification.
\item
  This note does not identify \(z_k\) or \(R\) with a curvature radius,
  projection radius, or internal radius.
\item
  This note does not identify the coordinate labels
  \(x,y,z_1,z_2,\dots\) with physical dimensions.
\item
  The ``imaginary-valued coordinate'' and ``invisible axis'' in this
  note are algebraic and display-level specifications of coordinate
  labels; they do not mean unobservable physical degrees of freedom or
  hidden physical dimensions.
\end{itemize}

These are boundary conditions that prevent the algebraic observation
from leaping to physical claims.

\begin{center}\rule{0.5\linewidth}{0.5pt}\end{center}

\section{7. Relation to the preceding three
notes}\label{relation-to-the-preceding-three-notes}

\subsection{7.1 Relation to the square-quantity-readout
paper}\label{relation-to-the-square-quantity-readout-paper}

Whereas the square-quantity-readout paper treated the individual
square-quantity readout \((x_i)\longmapsto(x_i^2)\), this note treats

\[
(x,y)\longmapsto Z=x+iy\longmapsto Z\bar Z=x^2+y^2 ,
\]

a two-component sum-of-squares readout by the conjugate complex norm,
and its equal-footing multi-dimensional sum
\(\sum_j Z_j\bar Z_j=\sum_n x_n^2\).

\subsection{7.2 Relation to the odd-harmonic isolated-peak
paper}\label{relation-to-the-odd-harmonic-isolated-peak-paper}

Whereas the second note treated ``what localized shapes can be made on a
phase interval,'' this note treats ``how a complex representation
carrying a phase angle connects to square-quantity readout.'' It targets
not a concrete odd-harmonic sum \(S_N(\varphi)\) but the phase
representation \(Z=\rho e^{i\theta}\) itself.

\subsection{7.3 Relation to the relative-phase contrast
paper}\label{relation-to-the-relative-phase-contrast-paper}

In the third note, giving two copies a common relative phase produced
\(\psi_\alpha=2\cos\alpha\,S_N\) and \(I_\alpha=4\cos^2\alpha\,I_N\), so
that the relative phase appeared as a squared-intensity coefficient
without changing the waveform. In this note, the conjugate product
\(Z\bar Z\) makes the phase vanish, leaving the norm square \(\rho^2\).
What the two share is that a real quantity appears \textbf{when a
phase-carrying structure is read as a square quantity}, not the phase
itself.

\begin{center}\rule{0.5\linewidth}{0.5pt}\end{center}

\section{8. Dangerous readings and safe
readings}\label{dangerous-readings-and-safe-readings}

\subsection{8.1 Dangerous readings (this note does not
assert)}\label{dangerous-readings-this-note-does-not-assert}

\begin{itemize}
\tightlist
\item
  Derived space or spacetime from complex numbers.
\item
  Phase space and real space are identical.
\item
  Derived special/general relativity from a complex norm.
\item
  The negative squared contribution is a spacetime metric signature.
\item
  The coordinate label \(z_k\) is time.
\item
  \(R\) or \(z_k\) is a curvature radius.
\item
  Coordinate labels are physical dimensions.
\end{itemize}

\subsection{8.2 Safe readings (what this note can
say)}\label{safe-readings-what-this-note-can-say}

\begin{enumerate}
\def\labelenumi{\arabic{enumi}.}
\tightlist
\item
  The complex polar form is a standard map sending a polar form
  containing a phase angle to two orthogonal components.
\item
  Its conjugate product is the sum of squares of two orthogonal real
  components, readable as a form of square-quantity readout.
\item
  Bundling real components two at a time generalizes to the
  equal-footing multi-dimensional sum of squares
  \(\sum_j Z_j\bar Z_j=\sum_n x_n^2=R^2\).
\item
  An invisible axis not displayed in the complex plane simply joins as
  an equal-footing axis with the same sign, and the radius \(R\) is
  unique.
\item
  The conjugate product \(X\bar X\) and the non-conjugate square \(X^2\)
  agree for real coordinates but diverge only for imaginary ones
  (\(X_h\bar X_h=z_h^2\) versus \(X_h^2=-z_h^2\)). The minus sign arises
  only from this divergence.
\item
  Expanding the zero-sum-of-squares form
  \(\sum_j Z_j\bar Z_j+\sum_k (i z_k)^2=0\)---the conjugate-norm
  backbone plus the imaginary-coordinate square---yields the algebraic
  relation \(\sum_n x_n^2=\sum_k z_k^2\), that the sum of squares of the
  real coordinates equals the sum of squares of the imaginary coordinate
  labels. Non-triviality requires at least one imaginary-valued
  coordinate.
\item
  All of this is a geometric observation on the algebraic arrangement of
  conjugate complex numbers and sums of squares, and makes no physical
  claim.
\end{enumerate}

\begin{center}\rule{0.5\linewidth}{0.5pt}\end{center}

\section{9. Conclusion}\label{conclusion}

In this note, as the fourth in the series following square-quantity
readout, odd-harmonic phase structure, and relative-phase contrast, we
observed the connection between the conjugate complex norm and
square-quantity readout, and generalized it to an equal-footing
multi-dimensional sum of squares.

For the complex polar form \(Z=\rho e^{i\theta}\), putting
\(x=\rho\cos\theta,\ y=\rho\sin\theta\) gives \(Z\bar Z=x^2+y^2=\rho^2\)
(a standard fact). The center of this note is to place this known norm
within the square-quantity-readout series and, further, to bundle real
components two at a time and generalize to the equal-footing
multi-dimensional sum of squares

\[
\sum_j Z_j\bar Z_j=\sum_n x_n^2=R^2 .
\]

There is no special axis here; an invisible axis also joins as an
equal-footing axis with the same sign, and \(R\) is unique.

Furthermore, adding the square of an imaginary-valued coordinate to this
conjugate-norm backbone, we observed the zero-sum-of-squares form

\[
\sum_j Z_j\bar Z_j+\sum_k (i z_k)^2=0 .
\]

Since the conjugate product \(X\bar X\) and the non-conjugate square
\(X^2\) agree for real coordinates and diverge only for imaginary ones
(\(X_h\bar X_h=z_h^2\) vs \(X_h^2=-z_h^2\)), expanding gives, through
\(i^2=-1\) for the purely imaginary-valued coordinate, the algebraic
relation

\[
\sum_{n} x_n^2=\sum_{k} z_k^2 .
\]

The negative squared contribution appearing here is a consequence of
\(i^2\), not a hand-placed metric signature, and no double definition of
\(R\) occurs. The conjugate-norm type \(Z\bar Z\) and the non-conjugate
square \(X_h^2\) of an imaginary coordinate are not the same operation,
and this note keeps the distinction.

The above records how the conjugate positive-definite norm
\(\sum_j Z_j\bar Z_j\) and the non-conjugate square term \((i z_k)^2\)
of an imaginary-valued coordinate can be connected within a single
zero-sum-of-squares form, together with consistency with square-quantity
readout and the equal-footing multi-dimensional sum of squares---an
observation limited to elementary algebra and geometric interpretation.
This note does not connect to physics (spacetime, relativity, quantum
theory, metric signature), and does not identify coordinate labels with
physical dimensions.

\begin{center}\rule{0.5\linewidth}{0.5pt}\end{center}

\section{References}\label{references}

\subsection{Self-references}\label{self-references}

{[}1{]} N. Kihara, ``Linear Simplexification by Square-Quantity
Readout,'' v0.1-r1, Zenodo DOI: 10.5281/zenodo.20785540, 2026.

{[}2{]} N. Kihara, ``An Observation on the Isolated Peak Wave of a
Constant-Amplitude Odd-Harmonic Sum on a Half-Wavelength Phase Interval
and Its Localization,'' v0.6, Zenodo DOI: 10.5281/zenodo.21073985, 2026.

{[}3{]} N. Kihara, ``Waveform Invariance and a Contrast Law for the
Two-Copy Common-Relative-Phase Superposition of a Half-Wavelength
Odd-Harmonic Isolated-Peak Wave,'' v0.1, Zenodo DOI:
10.5281/zenodo.20923462, 2026.

\subsection{Related standard material}\label{related-standard-material}

{[}4{]} T. Takagi, \emph{Kaiseki Gairon} (Introduction to Analysis), 3rd
revised ed., Iwanami Shoten, 1961.

{[}5{]} Standard textbook material on the polar form of complex numbers,
Euler's formula, the complex norm, and standard quadratic forms.

\begin{center}\rule{0.5\linewidth}{0.5pt}\end{center}

\section{Appendix A: Minimal chain of formulas
(revised)}\label{appendix-a-minimal-chain-of-formulas-revised}

\begin{enumerate}
\def\labelenumi{\arabic{enumi}.}
\tightlist
\item
  In the complex plane, set \(Z=x+iy\) (role (1) of \(i\): bundling).
\item
  By the conjugate product, \(Z\bar Z=x^2+y^2=\rho^2\).
\item
  Bundle real components two at a time and add on an equal footing:
  \(\displaystyle\sum_j Z_j\bar Z_j=\sum_n x_n^2=R^2\). An invisible
  axis merely joins as an equal-footing axis with the same sign.
\item
  Confirm the difference between conjugate product and square: for real
  coordinates \(x\bar x=x^2\); for imaginary coordinates they diverge,
  \(X_h\bar X_h=z_h^2\) versus \(X_h^2=(i z_h)^2=-z_h^2\) (role (2) of
  \(i\)). The minus sign arises only here.
\item
  The zero-sum-of-squares form adding the imaginary-coordinate square to
  the conjugate-norm backbone:
  \(\displaystyle\sum_j Z_j\bar Z_j+\sum_k (i z_k)^2=0\). Non-triviality
  requires an imaginary-valued coordinate.
\item
  Expand and rearrange:
  \(\displaystyle\sum_j Z_j\bar Z_j=\sum_k z_k^2\),
  i.e.~\(\displaystyle\sum_{n} x_n^2=\sum_{k} z_k^2\). The minus sign is
  a consequence of \(i^2\). The \(Z\bar Z\) type and \(X^2\) type are
  not confused, and no double definition of \(R\) occurs.
\end{enumerate}

\begin{center}\rule{0.5\linewidth}{0.5pt}\end{center}

\section{Appendix B: Candidate connections for a next paper (outside the
scope of this
note)}\label{appendix-b-candidate-connections-for-a-next-paper-outside-the-scope-of-this-note}

This note was limited to observations on algebraic arrangement.
Candidates that a next paper might examine (none treated here):

\begin{enumerate}
\def\labelenumi{\arabic{enumi}.}
\tightlist
\item
  Organizing the geometric motivation for introducing imaginary-valued
  coordinates (a criterion for which axis to place as imaginary).
\item
  The geometry of the solution set (a complex quadric hypersurface) of
  the zero-sum-of-squares \(\sum_n X_n^2=0\).
\item
  Projection onto the real subspace and a generalization of the role of
  the invisible axis.
\item
  Comparison with the projection radius of the central-projection series
  (comparison only; no identification).
\item
  Making explicit the additional assumptions required if one proceeds to
  a physical interpretation (this note does not proceed there).
\end{enumerate}

\begin{center}\rule{0.5\linewidth}{0.5pt}\end{center}

\section{Appendix C: Conjugate-pair closure and the entry point to
non-recoverability}\label{appendix-c-conjugate-pair-closure-and-the-entry-point-to-non-recoverability}

In §5, adding the non-conjugate square \((i z_k)^2=-z_k^2\) of
imaginary-valued coordinates to the conjugate-norm-type square quantity
\(\sum_j Z_j\bar Z_j\) gave the non-trivial zero-sum-of-squares
\(\sum_j Z_j\bar Z_j+\sum_k (i z_k)^2=0\). This appendix records a very
simple observation that appears when this construction is read in
reverse. \textbf{It does not derive an uncertainty principle.}

\subsection{C.1 With real conjugate pairs alone, the zero-sum closes
trivially}\label{c.1-with-real-conjugate-pairs-alone-the-zero-sum-closes-trivially}

If all components are real-valued and bundled with no remainder into
fixed pairs \(Z_j=q_{2j-1}+i\,q_{2j}\) (\(q_n\in\mathbb{R}\)), then

\[
\sum_j Z_j\bar Z_j=\sum_n q_n^2\ge0
\]

is positive-definite, and \(\sum_n q_n^2=0\) admits only the all-zero
solution \(q_1=\cdots=q_{2m}=0\). That is, in this range the zero-sum
\textbf{closes trivially} (has no non-trivial solution).

\subsection{C.2 A non-trivial zero-sum requires a non-real
component}\label{c.2-a-non-trivial-zero-sum-requires-a-non-real-component}

Therefore, to obtain a non-trivial zero-sum one needs a negative
contribution such as \((i z_k)^2=-z_k^2\), i.e.~\textbf{at least one
non-real component} (or a breaking of the fixed conjugate-pair
structure) (§5.3). Put differently, real conjugate pairs alone are
positive-definite and \textbf{close too tightly} to produce a
non-trivial zero.

\subsection{C.3 The conjugate-norm readout collapses the phase
(many-to-one)}\label{c.3-the-conjugate-norm-readout-collapses-the-phase-many-to-one}

On the other hand, the conjugate-norm readout itself drops the phase,
\textbf{regardless of whether the system is closed}. For
\(Z=\rho e^{i\theta}\) we have \(Z\bar Z=\rho^2\), and \(\theta\)
vanishes. Since every \(Z=\rho e^{i\theta}\) with the same \(\rho\)
returns the same \(\rho^2\), the map

\[
Z\longmapsto Z\bar Z=\rho^2
\]

is \textbf{many-to-one}. Hence the original phase representation \(Z\)
cannot be uniquely recovered from the square quantity \(\rho^2\).

\subsection{C.4 The entry point to
non-recoverability}\label{c.4-the-entry-point-to-non-recoverability}

Combining C.1--C.3, square-quantity readout has \textbf{two
complementary facets}.

\begin{itemize}
\tightlist
\item
  \begin{enumerate}
  \def\labelenumi{(\roman{enumi})}
  \tightlist
  \item
    \textbf{Triviality of the zero}: as long as it closes into real
    conjugate pairs, the zero-sum is positive-definite and trivial;
    making it non-trivial requires a non-real component (C.1, C.2).
  \end{enumerate}
\item
  \begin{enumerate}
  \def\labelenumi{(\roman{enumi})}
  \setcounter{enumi}{1}
  \tightlist
  \item
    \textbf{Many-to-one readout}: the conjugate-norm readout always
    identifies the phase direction, regardless of closure (C.3).
  \end{enumerate}
\end{itemize}

This note calls the many-to-one-ness of (ii) the \textbf{entry point to
non-recoverability}. This is not a hidden-variable claim that ``there is
a true hidden phase we cannot observe,'' but the observation that
\textbf{the readout map does not distinguish the phase direction} at all
------ a degeneracy. Note that the phase erasure does not arise from a
breaking of closure; it is intrinsic to the conjugate-norm readout.

\subsection{C.5 Positioning and
non-claims}\label{c.5-positioning-and-non-claims}

This appendix does not derive the uncertainty principle of quantum
mechanics, does not replace it, and is not a substitute for the
non-commutative operator formalism. That a complex quantum theory
emerges from putting a complex structure \(J^2=-1\) on a real Hilbert
space, that almost-complex structures are tied to even-dimensionality,
and that the Born rule \(\psi\mapsto|\psi|^2\) drops phase information,
are all known standard facts. All this appendix adds is a rereading:
making explicit, as an \textbf{entry point} for distinguishing the
determinacy and non-recoverability of a square-quantity readout, the
prior question of ``whether the components close as fixed conjugate
pairs.'' The stance of not identifying coordinate labels with physical
dimensions, and not identifying \(z_k\) with time or a curvature radius
(§0, §6), is retained in this appendix as well.
\end{document}
